% ============================================================
%  Résumé (FR) & Abstract (EN) — Mémoire de Master BIBDA
%  CHOUFA Zouhair
% ============================================================

% ─────────────────────────────────────────────────────────────
%  RÉSUMÉ (FRANÇAIS)
% ─────────────────────────────────────────────────────────────
\clearpage
% \thispagestyle{empty}

% Local adjustment to reduce top spacing and center the title
\begin{center}
    \vspace*{-1cm} % Pulls the title up into the empty white space
    {\Huge\bfseries Résumé}
    \vspace{0.5cm} % Space before the paragraph starts
\end{center}
\addcontentsline{toc}{chapter}{Résumé}

La dimension « Variété » du Big Data se traduit, en contexte
industriel, par la coexistence de multiples sources de données
tabulaires hétérogènes (formats, schémas, conventions de nommage),
rendant leur intégration sémantique difficile à automatiser à grande
échelle. Les Knowledge Graphs (KG), fondés sur le modèle RDF et les
standards du Web sémantique, offrent un cadre formel pour unifier ces
sources autour d'une ontologie commune ; toutefois, les pipelines de
construction existants restent le plus souvent fragmentés,
partiellement manuels, ou insuffisamment évalués de bout en bout.

Ce mémoire propose un pipeline complet, modulaire et reproductible de
transformation de données tabulaires hétérogènes en Knowledge Graph.
L'architecture proposée s'articule en quatre modules~: (i) ingestion
multi-source et profilage automatique de la qualité des données ;
(ii) \textit{schema matching} hybride combinant une approche
symbolique par embeddings Sentence-BERT et un agent LLM (LLaMA~3.1
via Groq, orchestré par LangChain), suivi de la génération
automatique de règles de mapping YARRRML ; (iii) triplification
sémantique déclarative des données au moyen du moteur Morph-KGC ;
(iv) résolution d'entités inter-sources par le modèle probabiliste de
Fellegi-Sunter (bibliothèque Splink), complétée par une validation
formelle du graphe produit via des contraintes SHACL. L'ensemble du
pipeline est conteneurisé avec Docker Compose afin de garantir sa
reproductibilité et son industrialisation potentielle.

L'approche a été évaluée sur un gold standard annoté, construit à
partir d'un jeu de données transactionnel augmenté artificiellement
(\texttt{Faker}) et scindé en deux sources hétérogènes simulant des
silos organisationnels distincts. Les résultats montrent que
l'architecture hybride de \textit{schema matching} atteint un
F1-score de 0,870 après calibration du seuil de décision
($\tau^{*} = 0{,}82$), égalant ou surpassant chacune de ses
composantes utilisée isolément tout en maîtrisant le coût des appels
LLM. Le module de triplification génère 82\,420 triplets RDF en
seulement 7 secondes, tandis que la résolution d'entités, évaluée sur
1\,600 paires réelles, atteint un F1-score de 0,990 après recalibrage
des règles de blocage ($\lambda^{*} = 0{,}90$). Une discussion
critique structurée autour des menaces classiques à la validité
expérimentale (interne, externe, de construit et de conclusion)
complète cette évaluation et identifie les limites du travail,
notamment le caractère synthétique du jeu de données et la
sensibilité du modèle de résolution d'entités à la configuration des
règles de blocage.

Ce travail démontre ainsi la faisabilité d'un pipeline hybride
symbolique-neuronal pour la construction automatisée et évaluable de
Knowledge Graphs à partir de données tabulaires hétérogènes, et ouvre
plusieurs perspectives~: validation sur des données réelles
anonymisées, définition automatique de clés de blocage adaptatives,
et substitution des appels LLM distants par des modèles open source
déployés localement afin de réduire la latence à grande échelle.

\vspace{0.8cm}
\noindent\textbf{Mots-clés~:} Big Data, Knowledge Graph, hétérogénéité
des données, \textit{schema matching}, Large Language Models,
RDF, YARRRML, Morph-KGC, résolution d'entités, SHACL, reproductibilité.

\newpage

% ─────────────────────────────────────────────────────────────
%  ABSTRACT (ENGLISH)
% ─────────────────────────────────────────────────────────────
% \thispagestyle{empty}

\begin{center}
    \vspace*{-1cm} % Pulls the Abstract up identically
    {\Huge\bfseries Abstract}
    \vspace{0.5cm}
\end{center}
\addcontentsline{toc}{chapter}{Abstract}

The "Variety" dimension of Big Data manifests, in industrial
settings, as the coexistence of multiple heterogeneous tabular data
sources differing in format, schema, and naming conventions, making
their semantic integration difficult to automate at scale. Knowledge
Graphs (KGs), built upon the RDF model and Semantic Web standards,
provide a formal framework for unifying such sources around a shared
ontology; however, existing construction pipelines remain, for the
most part, fragmented, partially manual, or insufficiently evaluated
end-to-end.

This thesis proposes a complete, modular, and reproducible pipeline
for transforming heterogeneous tabular data into a Knowledge Graph.
The proposed architecture is organized into four modules: (i)
multi-source ingestion and automated data quality profiling; (ii) a
hybrid schema matching approach combining a symbolic
Sentence-BERT embedding method with an LLM agent (LLaMA~3.1 via Groq,
orchestrated through LangChain), followed by the automatic generation
of YARRRML mapping rules; (iii) declarative semantic triplification
of the data using the Morph-KGC engine; (iv) cross-source entity
resolution based on the Fellegi-Sunter probabilistic model (Splink
library), complemented by a formal validation of the resulting graph
through SHACL constraints. The entire pipeline is containerized with
Docker Compose to ensure reproducibility and potential
industrialization.

The approach was evaluated on an annotated gold standard, built from
a transactional dataset artificially augmented via \texttt{Faker} and
split into two heterogeneous sources simulating distinct
organizational silos. Results show that the hybrid schema matching
architecture achieves an F1-score of 0.870 after calibrating the
decision threshold ($\tau^{*} = 0.82$), matching or outperforming
each of its components used in isolation while keeping LLM API costs
under control. The triplification module generates 82,420 RDF
triples in only 7 seconds, while entity resolution, evaluated on
1,600 real pairs, reaches an F1-score of 0.990 after recalibrating
the blocking rules ($\lambda^{*} = 0.90$). A critical discussion
structured around the classical threats to experimental validity
(internal, external, construct, and conclusion validity) complements
this evaluation and identifies the work's limitations, notably the
synthetic nature of the dataset and the sensitivity of the entity
resolution model to blocking rule configuration.

This work thus demonstrates the feasibility of a hybrid
symbolic-neural pipeline for the automated and evaluable construction
of Knowledge Graphs from heterogeneous tabular data, and opens
several avenues for future work: validation on real anonymized data,
automatic definition of adaptive blocking keys, and the substitution
of remote LLM calls with locally deployed open-source models to
reduce latency at scale.

\vspace{0.8cm}
\noindent\textbf{Keywords:} Big Data, Knowledge Graph, data
heterogeneity, schema matching, Large Language Models, RDF, YARRRML,
Morph-KGC, entity resolution, SHACL, reproducibility.

\newpage